\section{Introduction}\label{BmaGretlSec1}
We know from elementary statistical theory that linear regression
attempts to model the relationship between two or more variables by
fitting a simple linear equation to observed data. In the classical
approach, we usually rely on ordinary least squares (OLS) or maximum
likelihood (ML) estimates and popular model selection criteria, i.e.,
the Akaike information criterion (AIC) and the Bayesian information
criterion (BIC), to find the ``best'' model. A problem with this
approach arises when we have to select a ``good'' subset of variables
from a large set of regressors. When the number of possible exogenous
variables is $K$, the number of possible linear models is $2^K$.  If
we have, for example, $K=30$ possible regressors, the number of
possible combinations of variables to be included or not equals
1073741824. This means that it is very difficult, if not impossible,
to find the estimates for all combinations. Moreover,
\cite{RafteryMadiganHoeting1997} show that standard variable selection
procedures lead to different estimates and conflicting conclusions
about the main questions of interest.

Bayesian model averaging is a useful alternative to other variable
selection procedures, because it incorporates model uncertainty into
conclusions about the estimated parameters.  Bayesian model averaging
is a standard Bayesian solution to model uncertainty, where the
inference on parameters is based on a weighted average over all
possible models under consideration, rather than on one single
regression model. These weights are the Bayesian posterior
probabilities of the individual models.

There is a recent and growing literature on Bayesian model
averaging. Examples of applications of Bayesian model averaging can be
found in a number of works (see, for example,
\citealt{HoetingMadiganRafteryVolinsky1999}, \citealt{Steel2011} for a
recent overview).  Our software package for parameter estimation and
model comparison of linear regression models is based on
\cite{FernandezLeySteel2001a,FernandezLeySteel2001b} and
\cite{Koop2003}. We use the Markov chain Monte Carlo model composition
($\mbox{MC}^3$) sampling algorithm developed by \cite{MadiganYork1995}
to select a representative subset of models.

\cite{DoppelhoferWeeks2005,DoppelhoferWeeks2009} define a jointness
measure of dependence among explanatory variables that appear in the
linear regression models. We use this measure to identify whether two
variables are substitutes, complements or neither.  A similar
jointness measure was also proposed by \cite{LeySteel2007}.

We propose Bayesian model averaging package for
\pkg{gretl}, \citep{cottrell_luchcetti_gretl_manual} named \pkg{BMA}.
We can list several reasons why, in our opinion, it is important to
address this topic. \pkg{gretl} is
an increasingly popular, free, open-source software for econometric
analysis, both for students and academics. It is written in pure
\proglang{C} ``by econometricians for econometricians'' with the
international community gathered around two men: Allin Cottrell and
Riccardo ``Jack'' Lucchetti. Since 2009 the community has been meeting
every two years at the \pkg{gretl} conference (see
\url{https://gretlconference.org} for details). \pkg{gretl} is
licensed under the GPL-3 license (\url{http://www.gnu.org/copyleft/gpl.html})
and has been reviewed several times in various international journals.
The best known reviews were written by
\cite{RePEc:jae:japmet:v:18:y:2003:i:1:p:105-110},
\cite{RePEc:jae:japmet:v:22:y:2007:i:4:p:849-854} and
\cite{Rosenblad:2008:JSSOBK:v25s01}. Some brief information about
\pkg{gretl} was also presented in \cite{Lucchetti:2010:JSSOBK:v41i11} and
\cite{RePEc:wly:japmet:v:26:y:2011:i:5:p:880-885}.  Unlike most other
statistical software packages \pkg{gretl} has an easy-to-use graphical user
interface (GUI). Our software package is, therefore, a free and
easy-to-use tool for Bayesian model averaging.

The rest of the paper is organized as follows:
Section~\ref{BmaGretlSec2} briefly outlines Bayesian model averaging
for linear regression models with the $\mbox{MC}^3$ sampling algorithm and
jointness measures.  Section~\ref{BmaGretlSec3} provides an overview
of the \pkg{gretl} packages for Bayesian model averaging.
Section~\ref{BmaGretlSec4} includes an overview of Bayesian model
averaging software packages.  Section~\ref{BmaGretlSec5} presents
an empirical illustration. The final section concludes. 
